\begin{abstract}
Persistent controversy in televised dance competitions reflects a structural tension between judges’ technical evaluations and audience voting: strong performers may exit under transient vote swings, whereas popularity-dominant contestants can become effectively “locked in,” eroding perceived competitive legitimacy. Because vote totals are unobserved, audience support must be inferred before the institutional effects of rule design can be audited.

For Problem 1, we propose an \textbf{inverse reconstruction framework} that models weekly audience support as a simplex-valued share vector and samples the rule-feasible region jointly implied by published procedures and observed eliminations. Operable weeks are screened by a feasibility-based diagnostic, and the reconstructed latent votes are validated on elimination weeks with a binary-outcome Posterior Predictive Consistency (PPC) of $90.5\%$.

For Problem 2, a \textbf{same-week counterfactual design} fixes each week’s realized information and varies only the voting aggregation rule. Rank aggregation stabilizes outcomes by buffering score--vote gaps yet can obscure substantive technical separation; percent aggregation preserves intensity of support and better reflects audience intent but magnifies the risk of extreme popularity. Counterfactual simulations show that voting aggregation rules change the eliminated contestant in $26\%$ of historical weeks. Judge-save layers can mitigate anomalous weeks, but opaque activation invites perceptions of discretionary interference.

For Problem 3, a \textbf{three-channel structure} couples judges’ scores, conditional audience preference, and an explicit survival-risk channel via discrete-time hazard modeling, enabling professional-partner contributions to be decomposed into technical uplift and popularity spillover and translated into changes in elimination hazard.

For Problem 4, PAGS (Professional--Adaptive Guard System) replaces discretionary intervention with rule-based safeguards through adaptive vote compression, phase-dependent weighting to strengthen late-stage technical credibility, and a minimal-intervention guardrail triggered only by statistically intolerable technical collapse. Historical simulations indicate that \textbf{the professionalism index (P, measuring alignment with technical scoring) increases from $0.419$ to $0.497$}. In the Bobby Bones stress test, PAGS re-calibrates a near-lock-in baseline to an $85.6\%$ elimination probability, aligning survival likelihood more closely with technical merit. Robustness is supported by sampling convergence, prior perturbations with negligible drift in key signals, and cross-model agreement of driver rankings. Findings are distilled into a producer-facing memo to guide transparent evaluation and revision of future rules.
\end{abstract}

\begin{keywords}
bayesian inverse reconstruction; counterfactual simulation; discrete-time hazard model; voting aggregation rules; PAGS architecture; agency--meritocracy trade-off.
\end{keywords}
